\section{Using the lab GUI}
\label{sec:lab-gui}

KeyQuorum Lab runs this crate in a browser as a small sandboxed machine
(Section~\ref{sec:overview}): every button and every terminal line runs the
real \code{keyquorum} / \code{keyquorum-device} CLI, compiled to
WebAssembly, against an in-browser machine with its own files, SQLite
stores, and relay. This section documents the lab's own controls --- the
same guide is built into the lab itself as a \textbf{Help} tab, so nothing
here should ever describe the GUI differently from what is actually on the
page.

\begin{noteBox}
Everything the lab GUI does, it does by running a real command. The
Activity tab (Section~\ref{sec:lab-activity}) shows the exact command line
behind whatever you last did, and the Terminal tab (Section~\ref{sec:lab-terminal})
lets you type that same command yourself. There is no separate ``demo
mode'' behind the buttons.
\end{noteBox}

\subsection{Before you start: identity and drives}

The bar at the top of the page (\textbf{Active user}) shows who you are
acting as right now, and which mock USB drive their personal key slot lives
on. Clicking another name's chip switches to them, the way sitting down at
a different desk would. Most actions need that person's drive
\emph{inserted} first --- a mock drive behaves like a real one: its slots
are only usable while it is connected, and moving a slot onto someone
else's drive needs both drives inserted at once, matching the physical
requirement of moving a token between two real USB drives
(Section~\ref{sec:devices-and-custody}).

\subsection{Organization key tree}

The org-wide split tree, drawn as the dotted-label hierarchy it is
(Section~\ref{sec:hierarchical-principals}): \code{M}, \code{M.S},
\code{M.S.1}, and so on. You only ever see your own slice --- your
lineage, your descendants, siblings, and established-bridge peers --- the
same slice \code{keyquorum tree fetch} would return for your label. A leaf
tagged \emph{required} contributed a share the last time you unlocked a
file; \emph{satisfied} means that share was actually presented.

\begin{itemize}
  \item \textbf{Revoke key} bans a leaf's hardware key outright
  (\code{keyquorum revoke}); it drops that leaf's bindings and pairings but
  does not evict it or refresh survivor shares.
  \item \textbf{Reissue\ldots} replaces a leaf's hardware key onto an
  already-provisioned replacement token (\code{keyquorum reissue}).
  \item \textbf{Bridges} whitelists and establishes cross-branch links
  (\code{bridge allow} / \code{bridge add}) --- the same two-step process
  described in Section~\ref{sec:splitting-secrets}.
  \item \textbf{Tree restructure} walks through proposing a new public
  generation and collecting the parent's countersignature when one is
  required (Section~\ref{sec:authenticated-updates}).
\end{itemize}

\subsection{Mock USB devices}

Each card is one device container: a signed device descriptor plus one
sealed token per slot. \textbf{Insert} / \textbf{Eject} plug a drive in or
pull it out. \textbf{Move} asks for a destination drive (both drives must
be inserted) and relocates a slot's keypair onto it --- the point where two
people's tokens land in one physical container and start counting as a
single device toward any multi-device quorum
(Section~\ref{sec:devices-and-custody}). \textbf{Copy} instead leaves the
original slot active and seals a second copy elsewhere, using the same
passphrase that already unlocks it.

\subsection{File Explorer}

A Windows-Explorer-style view onto the active user's files, grouped into
folders. Double-clicking a file (or selecting it and pressing
\textbf{Open}) unlocks and views it --- a quorum file prompts for the
shares it needs, a password file for its password. \textbf{Send\ldots}
seals a copy to another label's key and drops it in their inbox;
\textbf{Sign} is only offered on public or received plaintext;
\textbf{Properties} shows the file's protection, size, and timestamps.
Right-clicking a row offers the same actions. A file with an expiry stamp
shows a status pill that flips to \emph{expired} on its own while the tab
is open --- the next unlock attempt (or a background scan) then deletes it
for real, the same destructive purge the CLI performs
(Section~\ref{sec:quorum-files}).

\subsection{Inbox and sent}

Letters are sealed to a recipient's key and routed by the mailbox relay
without it ever reading them --- sender and contents stay hidden until
\textbf{Receive} unseals a letter with your own slot (\code{keyquorum
deliver open}). \textbf{Reject} answers with a signed refusal instead. The
Sent view tracks whether a delivery is still awaiting the recipient's
signed acknowledgement, and \textbf{Check relay for acknowledgements} pulls
any that are waiting (Section~\ref{sec:mailbox-relay}).

\subsection{Security \& devices}

Everything that is not the hardware-key quorum tree or a file: locking a
note with a password (and optionally a PIN), exporting a file or credential
as a portable sealed bundle, minting a time-limited share link,
provisioning a fresh slot on an inserted drive and registering it as a new
leaf, signing or verifying with a private sign bridge you belong to
(Section~\ref{sec:private-bridges}), and, at the bottom, read-only device
logs and relay counters for the in-process mailbox this lab runs against.

\subsection{Activity and access trace}
\label{sec:lab-activity}

Every action taken leaves a trace here: which checks ran, in order, and
whether each one passed. The \textbf{Equivalent operation} line under the
latest trace is the literal command line that action ran --- the same one
that could be typed in the Terminal tab. The full log below keeps every
action for the session, not just the last one.

\subsection{Terminal}
\label{sec:lab-terminal}

A real shell on the same lab machine. Every button on every other tab runs
a command here under the hood; typing it directly works identically,
including commands no button exposes yet. A few starting points:

\begin{lstlisting}
keyquorum-device list /media/alice-usb
bridge list 1
usb insert david
su david
\end{lstlisting}

\code{usb insert} / \code{usb eject} and \code{su} are lab-only shell
verbs for drive and identity control; every other line is an ordinary
\code{keyquorum} / \code{keyquorum-device} command.
